\chapter{Medication (pregnancy)}

\section*{Definition and Overview}
Medication use during pregnancy requires a careful balance between the therapeutic benefits to the mother and the potential risks to the developing fetus. Many pregnant women require pharmacological treatment for pre-existing chronic conditions (such as asthma, epilepsy, diabetes, or depression) or acute gestational ailments (such as morning sickness, infections, or pre-eclampsia). The physiological changes of pregnancy, including expanded blood volume, increased glomerular filtration rate, and altered hepatic enzyme activity, significantly modify drug pharmacokinetics, often necessitating dose adjustments. Evaluating safety involves understanding teratogenicity (the ability to cause structural birth defects) and fetotoxicity (adverse effects on fetal growth or organ function). In Australia, prescribing is guided by the Therapeutic Goods Administration (TGA) prescribing categories, which classify drugs based on clinical safety data. Clinically, decision-making must be highly individualised, involving shared discussion of risks, benefits, and alternative therapies.

\section*{Epidemiology}
Pharmacological exposure during pregnancy is exceptionally common. Global epidemiological data indicate that over 80\% of pregnant women take at least one medication (over-the-counter or prescription) during pregnancy, with the average being 3 to 4 drugs. In Australia, the prevalence of medication use matches these figures, driven by increasing maternal age and higher rates of pre-existing chronic illnesses. The most frequently used therapeutic classes include vitamins and mineral supplements (such as folic acid and iron), antiemetics, analgesics (primarily paracetamol), and antibiotics. While most exposures involve clinically established safe agents, approximately 2\% to 3\% of congenital anomalies are estimated to be associated with maternal drug exposure, highlighting the clinical significance of structured risk evaluation and public health education.

\section*{Aetiology and Pathophysiology}
The potential for a medication to cause adverse fetal effects is determined by gestational timing, placental transfer, and drug characteristics.
\begin{itemize}
    \item \textbf{Gestational timing (first trimester):} The period of organogenesis (weeks 3 to 8 post-conception), during which the fetus is most susceptible to structural malformations (teratogenesis) caused by drugs.
    \item \textbf{Gestational timing (second and third trimesters):} A phase of histogenesis and functional development, where drug exposure can lead to growth restriction, organ dysfunction (fetotoxicity), or neonatal withdrawal.
    \item \textbf{Placental transfer:} Most drugs cross the placenta via passive diffusion; lipophilic, non-ionised, and low-molecular-weight substances cross more readily, exposing the fetal circulation.
    \item \textbf{Maternal pharmacokinetic alterations:} Increased plasma volume and renal clearance lower maternal drug concentrations, potentially compromising therapeutic efficacy unless doses are adjusted.
    \item \textbf{Fetal metabolic immaturity:} The fetal liver and kidneys have limited capacity to metabolise and excrete drugs, leading to prolonged exposure and accumulation in amniotic fluid.
\end{itemize}

\section*{Clinical Features}
The clinical features of medication exposure during pregnancy are evaluated in terms of maternal therapeutic response and potential neonatal syndromes.
\begin{itemize}
    \item \textbf{Maternal efficacy changes:} Decreased therapeutic control of chronic conditions (e.g. breakthrough seizures in epilepsy or loss of glycaemic control in diabetes) due to gestational clearance changes.
    \item \textbf{Structural anomalies:} Congenital defects (e.g. neural tube defects or cardiac anomalies) detected during routine fetal anomaly ultrasounds.
    \item \textbf{Fetal growth restriction:} Sub-optimal fetal growth and low birth weight, associated with chronic exposure to certain vasoactive or toxic drugs.
    \item \textbf{Neonatal abstinence syndrome (NAS):} Withdrawal symptoms in the newborn (including high-pitched crying, tremors, hypertonicity, and poor feeding) following maternal use of opioids, SSRIs, or benzodiazepines close to term.
    \item \textbf{Gestational complications:} Increased maternal risks, such as drug-induced gestational hypertension or alterations in uterine contractility.
\end{itemize}

\section*{Differential Diagnosis}
When evaluating fetal anomalies or neonatal distress, clinicians must distinguish drug-induced aetiologies from genetic, environmental, and infectious causes.
\begin{itemize}
    \item \textbf{Genetic or chromosomal anomalies:} Aneuploidies (e.g. Down syndrome) or microdeletions presenting with structural defects that mimic drug-induced teratogenesis.
    \item \textbf{Congenital TORCH infections:} Maternal infections (Toxoplasmosis, Rubella, Cytomegalovirus, Herpes Simplex) causing microcephaly, calcifications, or growth restriction that mimic drug fetotoxicity.
    \item \textbf{Maternal disease state effects:} Untreated or poorly controlled maternal illness (e.g. diabetic embryopathy or maternal thyroid dysfunction) causing birth defects rather than the prescribed medications.
    \item \textbf{Nutritional deficiencies:} Inadequate maternal intake (e.g. folate deficiency leading to neural tube defects) presenting similarly to antifolate drug exposure.
    \item \textbf{Idiopathic congenital defects:} Spontaneous malformations occurring in approximately 2--3\% of all births without a clear genetic, infectious, or pharmacological cause.
\end{itemize}

\section*{When to Seek Medical Care}
Pregnant women or those planning pregnancy should consult their general practitioner, obstetrician, or pharmacist before starting, changing, or discontinuing any medication. Urgent medical attention must be sought for acute complications or signs of drug toxicity.

\textbf{Call 000 (or local emergency services) for:}
\begin{itemize}
    \item \textbf{Acute drug toxicity or overdose:} Accidental or intentional ingestion of unsafe doses, presenting with altered consciousness, seizures, or respiratory depression.
    \item \textbf{Severe allergic reaction (anaphylaxis):} Difficulty breathing, facial or tongue swelling, or widespread hives following a new medication.
    \item \textbf{Status epilepticus:} Continuous or repetitive seizures in a pregnant patient, posing an immediate threat to maternal and fetal oxygenation.
    \item \textbf{Hypertensive crisis:} Severe headache, vision changes, or upper abdominal pain, indicating severe pre-eclampsia or eclampsia.
\end{itemize}

\section*{Diagnosis and Investigations}
Assessment involves verifying drug exposure history, monitoring maternal drug levels, and evaluating fetal well-being.
\begin{itemize}
    \item \textbf{Comprehensive medication history:} Documenting all prescription, over-the-counter, herbal, and recreational substances taken, including dosage and gestational timing.
    \item \textbf{Therapeutic drug monitoring (TDM):} Serial measurements of maternal blood levels for narrow-therapeutic-index drugs (e.g. anti-epileptics like lamotrigine) to guide dose adjustments.
    \item \textbf{High-resolution fetal ultrasound:} Performing morphology scans at 18--20 weeks to detect structural birth defects and monitor fetal growth and liquor volume.
    \item \textbf{Fetal echocardiography:} Specialised ultrasound of the fetal heart, indicated if the mother was exposed to potential cardiotoxic agents (e.g. lithium).
    \item \textbf{Neonatal assessment:} Documenting APGAR scores, physical examination, and monitoring for withdrawal symptoms or metabolic disturbances postpartum.
\end{itemize}

\section*{Management}
The management of medications in pregnancy focuses on optimising maternal health while minimising fetal risk through evidence-based selection.
\begin{itemize}
    \item \textbf{Use of established safe alternatives:} Selecting medications with extensive pregnancy safety records (e.g. using methyldopa or labetalol for hypertension; paracetamol for pain).
    \item \textbf{Pre-conception planning:} Optimising drug regimens prior to pregnancy, switching to safer drugs, and initiating high-dose folic acid (5 mg/day) for high-risk patients.
    \item \textbf{Dosing optimisation:} Splitting doses or adjusting quantities to maintain therapeutic maternal levels while avoiding high peak drug concentrations.
    \item \textbf{Non-pharmacological measures:} Emphasising lifestyle changes, physical therapy, or cognitive behavioural therapy to manage conditions like mild depression or back pain.
    \item \textbf{Neonatal surveillance and support:} Planning delivery at a centre equipped with a neonatal intensive care unit (NICU) when high-risk medication exposure occurs close to term.
\end{itemize}

\section*{Complications}
\begin{itemize}
    \item \textbf{Congenital malformations:} Major structural defects (e.g. cleft lip or palate, spina bifida, or cardiac septal defects) resulting from first-trimester exposure to teratogens.
    \item \textbf{Neonatal withdrawal (abstinence syndrome):} Irritability, feeding difficulties, and respiratory distress in the newborn following delivery.
    \item \textbf{Preterm birth or low birth weight:} Increased risk of premature delivery or intrauterine growth restriction associated with specific drug classes.
    \item \textbf{Long-term neurodevelopmental effects:} Cognitive or behavioural issues in children exposed in utero to certain antiepileptics (e.g. sodium valproate).
    \item \textbf{Maternal therapeutic failure:} Under-treatment of maternal illness due to fear of fetal risk, leading to severe maternal morbidity.
\end{itemize}

\section*{Prognosis and Outcomes}
The prognosis for pregnancy outcomes following medication exposure is highly favourable for the vast majority of commonly used drugs. When medications are selected carefully, chronic maternal conditions can be managed effectively without adverse fetal effects, leading to healthy term deliveries. In cases of exposure to known teratogens (such as sodium valproate or retinoids), the risk of adverse outcomes increases significantly, necessitating detailed genetic and fetal counselling. Establishing a multidisciplinary care team (consisting of the obstetrician, GP, pharmacist, and relevant specialists) is the single most important factor in optimising both maternal health and fetal prognosis.

\section*{Prevention and Self-Care}
\begin{itemize}
    \item \textbf{Consult before conceiving:} Discussing pregnancy plans with a GP to review and adjust chronic medications.
    \item \textbf{Avoid self-medication:} Refraining from taking over-the-counter drugs or herbal supplements without consulting a healthcare professional.
    \item \textbf{Take recommended supplements:} Consuming daily folic acid and iodine as recommended by Australian health guidelines.
    \item \textbf{Read medicine labels carefully:} Checking TGA pregnancy categories and patient information leaflets.
    \item \textbf{Maintain a health diary:} Keeping a log of all medications, doses, and dates taken during pregnancy.
\end{itemize}

\section*{Sources and Further Reading}
The following reputable organisations provide evidence-based databases and resources on medication safety during pregnancy:
\begin{itemize}
    \item MotherSafe NSW. \textit{Information service on medications, chemicals, and radiation in pregnancy.} https://www.sydneychildrenshospital.network/
    \item Therapeutic Goods Administration. \textit{Australian categorisation system for prescribing medicines in pregnancy.} https://www.tga.gov.au/
    \item Royal Women's Hospital. \textit{Pregnancy and breastfeeding medicines guide.} https://www.thewomens.org.au/
    \item Mayo Clinic. \textit{Medication during pregnancy: what's safe and what's not.} https://www.mayoclinic.org/
    \item UK Teratology Information Service. \textit{Evidence-based summaries on fetal drug exposure.} https://www.uktis.org/
\end{itemize}
